\documentclass[11pt]{article}

\usepackage{amsmath}
\usepackage{amsthm}
\usepackage{amssymb}

\theoremstyle{plain}
\newtheorem{theorem}{Theorem}
\newtheorem{lemma}{Lemma}
\newtheorem{proposition}{Proposition}
\newtheorem{corollary}{Corollary}
\theoremstyle{remark}
\newtheorem{remark}{Remark}

\newcommand{\SYT}{\operatorname{SYT}}
\newcommand{\cont}{\operatorname{content}}
\newcommand{\mincost}{\mathsf{mincost}}
\newcommand{\des}{\operatorname{des}}
\newcommand{\Trs}{T_{\mathrm{rs}}}

\title{Toward the Core Lemma:\\ the all-stay lower bound, and the swap-back obstruction}
\author{Clio}
\date{2026-05-22}

\begin{document}
\maketitle

\section{Introduction and context}

This note accompanies the companion paper
\texttt{2026-05-21-omega-trace-\linebreak[1]reciprocity-mindegree.tex}, in
which the identity
\[
  \min\deg_q \chi^\lambda(\Omega) \;=\; n(\lambda)
\]
for the Hecke-algebra character of the operator $\Omega$ is \emph{reduced}
to a purely combinatorial statement, the \emph{Core Lemma}: a minimum-cost
closed walk on standard Young tableaux $\SYT(\lambda)$, taken along the
staircase word, has total cost $\ge n(\lambda)$. The matching upper bound
is already established in the companion paper (an explicit all-stay walk at
the row-superstandard tableau realises cost $n(\lambda)$), so the Core Lemma
is exactly the missing lower bound.

The present note records partial progress toward the Core Lemma. Its
contributions are threefold:
\begin{itemize}
  \item a computational certificate that $\mincost(\lambda)=n(\lambda)$ for
        all feasible shapes through $n\le 7$ (\S\ref{sec:cert});
  \item a complete, unconditional proof of the \emph{all-stay lower bound}
        (Lemma~\ref{lem:allstay}, \S\ref{sec:allstay}), which closes the
        Core Lemma on the sub-class of walks that never swap, and which
        corrects the charge geometry sketched in \S7 of the companion paper;
  \item a precise diagnosis of why the natural dual potential fails on
        swap edges --- the \emph{swap-back obstruction} (\S\ref{sec:obstruction})
        --- which isolates the closed-walk constraint as the structure any
        full proof must exploit.
\end{itemize}
The Core Lemma itself (the lower bound over \emph{all} closed walks) remains
open; \S\ref{sec:status} states the status plainly and proposes a route.

\section{Setup (recap)}\label{sec:setup}

Let $\lambda\vdash n$. We index rows and columns from $0$. The
\emph{content} of cell $(r,c)$ is $\cont(r,c)=c-r$. For $T\in\SYT(\lambda)$
write $\cont_T(v)$ for the content of the cell holding value $v$.

For $i\in\{1,\dots,n-1\}$ set
\[
  d_T(i)\;=\;\cont_T(i{+}1)-\cont_T(i),
\]
the \emph{axial distance} between consecutive values, and define the
\emph{content-descent indicator}
\[
  \delta_T(i)\;=\;
  \begin{cases}
    1 & \text{if } \cont_T(i{+}1)<\cont_T(i)\ \ (\text{i.e. } d_T(i)<0),\\
    0 & \text{otherwise.}
  \end{cases}
\]

\paragraph{The staircase word.}
Let $\mathbf{w}_0$ be the word with letters $i_1,\dots,i_L$,
$L=\binom{n}{2}$, given by the concatenation of blocks
\[
  (1)\,(2\,1)\,(3\,2\,1)\,\cdots\,(n{-}1\ \cdots\ 1).
\]
Letter $i$ occurs exactly $n-i$ times in $\mathbf{w}_0$.

\paragraph{Closed walks and cost.}
A closed walk starts at a tableau $T=U_0\in\SYT(\lambda)$ and proceeds
$U_{k-1}\to U_k$, one step per letter $i_k$ of $\mathbf{w}_0$, with two
move types:
\begin{itemize}
  \item \textbf{STAY:} $U_k=U_{k-1}$;
  \item \textbf{SWAP:} $U_k=s_{i_k}U_{k-1}$, permitted only when the axial
        distance satisfies $|d_{U_{k-1}}(i_k)|\ge 2$.
\end{itemize}
The walk is \emph{closed} when $U_L=U_0$. The \emph{cost} of step $k$ is
$\delta_{U_k}(i_k)$, i.e.\ it is $1$ exactly when $i_k{+}1$ has strictly
lower content than $i_k$ in the \emph{after}-state $U_k$. The same-column
case $d=-1$ is forbidden: it can never be a visited state (a swap there is
disallowed, and it cannot arise in a valid $\SYT$ as an after-state of a
permitted move). We write
\[
  \mincost(\lambda)\;=\;\min_{\text{closed walks}}\ \sum_{k=1}^{L}\delta_{U_k}(i_k).
\]

\paragraph{The target quantity.}
With $0$-indexed rows,
\[
  n(\lambda)\;=\;\sum_{r\ge 0} r\,\lambda_r
            \;=\;\sum_{i\ge 1}(i-1)\lambda_i
            \;=\;\sum_{j}\binom{\lambda'_j}{2},
\]
where $\lambda'$ is the conjugate partition. The Core Lemma asserts
$\mincost(\lambda)\ge n(\lambda)$.

\section{Computational certificate}\label{sec:cert}

The minimum-cost closed walk is a shortest-path problem on the layered
state graph whose layers are indexed by the $L+1$ positions of
$\mathbf{w}_0$ and whose nodes in each layer are the tableaux of
$\SYT(\lambda)$, with STAY/SWAP edges as in \S\ref{sec:setup}. Solving the
LP dual (a longest-path / node-potential feasibility problem) yields, in
every case below, an \emph{integral} node potential certifying the lower
bound exactly.

\begin{proposition}[computational]\label{prop:cert}
For every $\lambda\vdash n$ with $n\le 7$ that is \emph{feasible} we have
\[
  \mincost(\lambda)\;=\;n(\lambda),
\]
and the dual LP admits an integral optimal potential. Here a shape is
infeasible, and excluded, precisely when it has two or more parts equal to
$1$: the staircase then forces a visited state with axial distance $d=-1$
(a same-column adjacency), which is forbidden. In all certified cases the
optimal walk is the all-stay walk at the row-superstandard tableau $\Trs$.
\end{proposition}

This is the empirical anchor for the conjectured Core Lemma. The proof of
Lemma~\ref{lem:allstay} below explains, unconditionally, why no \emph{all-stay}
walk can beat $n(\lambda)$; the residual gap is exactly the swap behaviour
analysed in \S\ref{sec:obstruction}.

\section{The all-stay lower bound}\label{sec:allstay}

For a fixed tableau $T$, the all-stay walk (no swaps at any step) has total
cost
\[
  D(T)\;:=\;\sum_{k=1}^{L}\delta_T(i_k)\;=\;\sum_{i=1}^{n-1}(n-i)\,\delta_T(i),
\]
since letter $i$ occurs $n-i$ times and the state never changes.

\begin{lemma}[All-stay lower bound]\label{lem:allstay}
For every $T\in\SYT(\lambda)$,
\[
  D(T)\;\ge\;n(\lambda),
\]
with equality at $T=\Trs$. Consequently no all-stay closed walk costs less
than $n(\lambda)$.
\end{lemma}

\begin{proof}
\textbf{Step 1 (reindex by threshold).}
Each weight $n-i$ counts the values strictly above $i$, namely
$n-i=\#\{v: v>i\}$. Hence
\[
  D(T)
   =\sum_{i:\,\delta_T(i)=1}(n-i)
   =\sum_{i:\,\delta_T(i)=1}\#\{v:v>i\}
   =\sum_{v=2}^{n}\des_{<v}(T),
\]
where
\[
  \des_{<v}(T):=\#\{\,i<v:\delta_T(i)=1\,\}
\]
is the number of content-descents at letters below $v$. (We may include
$v=1$ on the right since $\des_{<1}(T)=0$.)

\medskip
\textbf{Step 2 (a column tower forces descents).}
Fix a value $v$, occupying cell $(\rho,c)$ with $\rho$ its $0$-indexed row
and $c$ its column. Consider column $c$ from the top down to $v$: the cells
\[
  (0,c),(1,c),\dots,(\rho,c)
\]
hold values $a_0<a_1<\dots<a_\rho=v$, strictly increasing by column-strictness
of standard tableaux. For each $t\in\{0,\dots,\rho-1\}$, the values $a_t$ and
$a_{t+1}$ sit in the column-adjacent cells $(t,c)$ and $(t{+}1,c)$, so
\[
  \cont_T(a_{t+1})-\cont_T(a_t)
   =\cont(t{+}1,c)-\cont(t,c)
   =-1.
\]
Thus over the value interval $[a_t,a_{t+1}]$ the net content change is
$-1<0$. A content-trajectory with no down-step (no $i$ with $\delta_T(i)=1$)
is non-decreasing; since the net change is negative there must be at least
one content-descent at some letter $i$ with $a_t\le i\le a_{t+1}-1$, i.e.\ in
the interval $[a_t,\,a_{t+1}-1]$.

\medskip
\textbf{Step 3 (disjointness and a per-value bound).}
The integer intervals $[a_t,a_{t+1}-1]$ for $t=0,\dots,\rho-1$ are pairwise
disjoint (they tile $[a_0,a_\rho-1]$) and all lie in $[1,v-1]$ because
$a_\rho=v$ and $a_0\ge 1$. Each contributes at least one descent at a letter
$i<v$, and these descents lie in disjoint letter-ranges, so they are
distinct. Therefore
\[
  \des_{<v}(T)\;\ge\;\rho\;=\;\rho_T(v),
\]
where $\rho_T(v)$ denotes the ($0$-indexed) row of the cell holding $v$.

\medskip
\textbf{Step 4 (sum over values).}
Summing the per-value bound over all $v$,
\[
  D(T)=\sum_{v=1}^{n}\des_{<v}(T)\;\ge\;\sum_{v=1}^{n}\rho_T(v)
       =\sum_{r\ge 0} r\,\lambda_r
       =n(\lambda),
\]
since the multiset of row-indices over all cells is determined by the shape
alone: row $r$ contains exactly $\lambda_r$ cells, each contributing $r$.

\medskip
\textbf{Equality at $\Trs$.}
For the row-superstandard tableau $\Trs$ (rows filled left-to-right,
top-to-bottom), content increases by $+1$ within a row and drops when
passing from the end of row $r$ to the start of row $r+1$. Thus
content-descents occur exactly at the row-end letters $i=P_r$
($P_r=\lambda_0+\cdots+\lambda_r$ being the position of the last value of
row $r$). For such a $T$, each per-value bound in Step~3 is met with
equality (the only descents below $v$ are the row-ends strictly above $v$'s
cell), and the weights telescope to $\sum_r r\,\lambda_r=n(\lambda)$.
Hence $D(\Trs)=n(\lambda)$.

Finally, an all-stay closed walk has total cost $D(U_0)\ge n(\lambda)$ for
its starting tableau $U_0$, which proves the last claim.
\end{proof}

\begin{remark}
This is precisely the ``column-pair forces a descent'' charge, but the cost
is charged to \emph{anti-diagonal} (strictly-lower-content) descents, not to
same-column adjacent steps. This corrects the heuristic charge sketch in \S7
of the companion paper, which attempted to inject the cost into same-column
adjacencies. In fact a cost-$1$ step has $i$ and $i{+}1$ landing on an
\emph{anti-diagonal}: e.g.\ consecutive values in cells $(1,2)$ and $(2,1)$,
where $\cont$ drops by $1$ but the cells are not in the same column. The
same-column configuration $d=-1$ is forbidden and never visited, so it can
carry no charge.
\end{remark}

\section{Why naive potentials fail: the swap-back obstruction}\label{sec:obstruction}

\paragraph{The dual potential method.}
A \emph{node potential} is a family $\phi_k(T)$, $0\le k\le L$, $T\in\SYT(\lambda)$,
such that on every edge of the layered graph (a step $k$ taking tail-state
$T$ to head-state $T'$) one has
\[
  \phi_k(T')-\phi_{k-1}(T)\;\le\;\mathrm{cost}(\text{edge}).
\]
Telescoping along any closed walk gives
\[
  \sum_{k=1}^{L}\mathrm{cost}\;\ge\;\phi_L(U_L)-\phi_0(U_0)
                                  =\phi_L(U_0)-\phi_0(U_0)
                                  \ge \min_T\bigl(\phi_L(T)-\phi_0(T)\bigr),
\]
using $U_L=U_0$. So any valid potential yields the lower bound
$\mincost\ge\min_T(\phi_L(T)-\phi_0(T))$.

\paragraph{The natural candidate.}
Write $c_k(T)=\sum_{j\le k}\delta_T(i_j)$ for the prefix all-stay cost of the
fixed tableau $T$ through step $k$ (so $c_0(T)=0$, $c_L(T)=D(T)$). Set
\[
  \phi_k(T)\;=\;c_k(T)-D(T)+n(\lambda).
\]
On a STAY edge $T\to T$ at step $k$, $\phi_k(T)-\phi_{k-1}(T)=\delta_T(i_k)$,
the exact cost --- so STAY edges are tight. And
\[
  \phi_L(T)-\phi_0(T)=\bigl(c_L(T)-D(T)+n(\lambda)\bigr)-\bigl(0-D(T)+n(\lambda)\bigr)=c_L(T)-c_0(T)
\]
gives terminal value $D(T)$; combined with Lemma~\ref{lem:allstay} this
would deliver $\min_T D(T)=n(\lambda)$ \emph{if the candidate were valid}.
It is not.

\begin{proposition}\label{prop:obstruction}
The candidate $\phi_k(T)=c_k(T)-D(T)+n(\lambda)$ is \emph{not} a valid node
potential: it satisfies all STAY-edge constraints with equality but violates
the SWAP-edge constraints exactly on shapes admitting a $D$-decreasing swap.
The number of violated swap edges grows like $n-3$ across the certified
range (e.g.\ $63$ of $546$ swap edges for $n\le 6$). Concretely, take
$\lambda=(2,2)$, step $k=2$, letter $i=2$, with
\[
  U=\begin{smallmatrix}1&3\\2&4\end{smallmatrix},
  \qquad
  V=s_2U=\begin{smallmatrix}1&2\\3&4\end{smallmatrix}.
\]
Then $D(U)=4$, $D(V)=2$, and the swap edge $U\to V$ at step $2$ gives
\[
  \phi_2(V)-\phi_1(U)\;=\;2\;>\;1\;=\;\delta_V(2),
\]
violating the edge constraint.
\end{proposition}

\paragraph{Diagnosis: the swap-back obstruction.}
A swap at letter $i$ costs $0$ iff it \emph{resolves} a descent --- a
``sorting'' swap, whose before-state had $\delta_U(i)=1$ (axial distance
$d\le-2$, hence $|d|\ge2$ permits the swap and the after-state has the two
values ascending in content) --- and costs $1$ iff it \emph{creates} one ---
a ``de-sorting'' swap (before-state ascending, after-state a descent). A
sorting swap lowers the future stay-cost $D$, so the candidate over-credits
it: it books the entire $D$-saving at the single swap edge, which is exactly
the violation above.

But in a \emph{closed} walk the performed swaps compose to the identity ---
we return to $U_0=U_L$. So every sorting swap must eventually be undone by a
de-sorting swap. The cost-$1$ de-sorting swaps pay back exactly what the
cost-$0$ sorting swaps saved. A correct potential must \emph{amortize} the
$D$-difference across the whole walk rather than booking it at a single
swap; equivalently, the proof must exploit the closed-walk constraint (the
performed swaps multiply to the identity), which the all-stay argument of
\S\ref{sec:allstay} never uses.

\section{Status and next step}\label{sec:status}

To state the situation plainly:
\begin{itemize}
  \item The \textbf{Core Lemma} --- the lower bound
        $\mincost(\lambda)\ge n(\lambda)$ over \emph{all} closed walks ---
        remains \textbf{open}.
  \item Lemma~\ref{lem:allstay} closes the \textbf{all-stay} case
        unconditionally: any walk that never swaps costs $\ge n(\lambda)$.
  \item Proposition~\ref{prop:obstruction} \textbf{isolates the obstruction}:
        the natural potential fails precisely at $D$-decreasing (sorting)
        swaps, and the failure is exactly compensated, in any closed walk,
        by the de-sorting swaps that restore $U_0$.
\end{itemize}

\paragraph{Proposed route.}
Read the walk as a wiring diagram in the spirit of the $0$-Hecke monoid:
STAY steps act as idempotents and SWAP steps as crossings. The closed
condition forces the crossings to pair up (the composite permutation is the
identity), so each sorting crossing is matched to a de-sorting crossing
further along the word. Quantifying this sorting/de-sorting balance --- a
matching between cost-$0$ and cost-$1$ swaps along the staircase --- should
amortize the $D$-difference and convert the all-stay bound of
\S\ref{sec:allstay} into the full Core Lemma.

\end{document}
